% from ega2-1-fr.tex
\oldpage[II]{71}

fini $\mathcal{E}_i$ de $\mathcal{S}_n$ tel que
$\mathcal{E}_i|U_i=(E_i)^{\sim}$
(\hyperref[I.9.4.7-fr]{I, 9.4.7}). Il est clair que le
sous-$\mathcal{O}_Y$-Module $\mathcal{E}$ de $\mathcal{S}_n$, somme des
$\mathcal{E}_i$, répond à la question (chaque section $f_{ij}$ étant telle
que pour tout $x\in X_{f_{ij}}$, il y a un voisinage affine
$V\subset X_{f_{ij}}$ de $x$ tel que $f|V$ soit une base de
$\Gamma(V,\mathcal{L}^{\otimes n})$).
\end{proof}

\section{Fibrés projectifs. Faisceaux amples}
\label{section:II.4-fr}

\subsection{Définition des fibrés projectifs.}
\label{subsection:II.4.1-fr}

\begin{definition}[4.1.1]
\label{II.4.1.1-fr}
Soient $Y$ un préschéma, $\mathcal{E}$ un $\mathcal{O}_Y$-Module
quasi-cohérent, $\mathbf{S}_{\mathcal{O}_Y}(\mathcal{E})$ la
$\mathcal{O}_Y$-Algèbre symétrique de $\mathcal{E}$
(\hyperref[II.1.7.4-fr]{1.7.4}) qui est quasi-cohérente
(\hyperref[II.1.7.7-fr]{1.7.7}). On appelle \emph{fibré projectif sur $Y$
défini par $\mathcal{E}$} et on note $\mathbf{P}(\mathcal{E})$ le $Y$-schéma
$P=\operatorname{Proj}(\mathbf{S}_{\mathcal{O}_Y}(\mathcal{E}))$. Le
$\mathcal{O}_P$-Module $\mathcal{O}_P(1)$ est appelé le \emph{faisceau
fondamental sur $P$}.
\end{definition}

Lorsque $Y$ est affine d'anneau $A$, on a $\mathcal{E}=\widetilde{E}$, où
$E$ est un $A$-module, et on écrit alors $\mathbf{P}(E)$ au lieu de
$\mathbf{P}(\widetilde{E})$.

Lorsque l'on prend $\mathcal{E}=\mathcal{O}_Y^n$, on écrit
$\mathbf{P}_Y^{n-1}$ au lieu de $\mathbf{P}(\mathcal{E})$~; si de plus $Y$
est affine d'anneau $A$, on écrit aussi $\mathbf{P}_A^{n-1}$ au lieu de
$\mathbf{P}_Y^{n-1}$. Comme
$\mathbf{S}_{\mathcal{O}_Y}(\mathcal{O}_Y)$ s'identifie canoniquement à
$\mathcal{O}_Y[T]$ (\hyperref[II.1.7.4-fr]{1.7.4}),
$\mathbf{P}_Y^0$ s'identifie canoniquement à $Y$
(\hyperref[II.3.1.7-fr]{3.1.7})~; l'exemple
(\hyperref[II.2.4.3-fr]{2.4.3}) n'est autre que $\mathbf{P}_K^1$.

\begin{env}[4.1.2]
\label{II.4.1.2-fr}
Soient $\mathcal{E}$, $\mathcal{F}$ deux $\mathcal{O}_Y$-Modules
quasi-cohérents~; $u:\mathcal{E}\to\mathcal{F}$ un
$\mathcal{O}_Y$-homomorphisme~; il lui correspond canoniquement un
homomorphisme
$\mathbf{S}(u):\mathbf{S}_{\mathcal{O}_Y}(\mathcal{E})\to
\mathbf{S}_{\mathcal{O}_Y}(\mathcal{F})$ de $\mathcal{O}_Y$-Algèbres
graduées (\hyperref[II.1.7.4-fr]{1.7.4}). Si $u$ est \emph{surjectif}, il en
est de même de $\mathbf{S}(u)$, et par suite
(\hyperref[II.3.6.2-fr]{3.6.2, (i)}),
$\operatorname{Proj}(\mathbf{S}(u))$ est une immersion fermée
$\mathbf{P}(\mathcal{F})\to\mathbf{P}(\mathcal{E})$, que nous noterons
$\mathbf{P}(u)$. On peut donc dire que $\mathbf{P}(\mathcal{E})$ est un
\emph{foncteur contravariant en $\mathcal{E}$}, à condition de ne prendre
pour morphismes des $\mathcal{O}_Y$-Modules quasi-cohérents que les
homomorphismes surjectifs.

En supposant toujours $u$ surjectif et posant
$P=\mathbf{P}(\mathcal{E})$, $Q=\mathbf{P}(\mathcal{F})$ et
$j=\mathbf{P}(u)$, on a, à un isomorphisme près,
\[
\label{II.4.1.2.1-fr}
  j^*(\mathcal{O}_P(n))=\mathcal{O}_Q(n)
  \qquad\text{pour tout }n\in\mathbf{Z}
\tag{4.1.2.1}
\]
comme il résulte de (\hyperref[II.3.6.3-fr]{3.6.3}).
\end{env}

\begin{env}[4.1.3]
\label{II.4.1.3-fr}
Soit maintenant $\psi:Y'\to Y$ un morphisme, et soit
$\mathcal{E}'=\psi^*(\mathcal{E})$~; on a alors
$\mathbf{S}_{\mathcal{O}_{Y'}}(\mathcal{E}')=
\psi^*(\mathbf{S}_{\mathcal{O}_Y}(\mathcal{E}))$
(\hyperref[II.1.7.5-fr]{1.7.5})~; par suite
(\hyperref[II.3.5.3-fr]{3.5.3}), on a
\[
\label{II.4.1.3.1-fr}
  \mathbf{P}(\psi^*(\mathcal{E}))=
  \mathbf{P}(\mathcal{E})\times_YY'
\tag{4.1.3.1}
\]
à un isomorphisme canonique près~; en outre, on a évidemment
\[
  \psi^*((\mathbf{S}_{\mathcal{O}_Y}(\mathcal{E}))(n))
  =(\mathbf{S}_{\mathcal{O}_{Y'}}(\mathcal{E}'))(n)
\]
pour tout $n\in\mathbf{Z}$, donc, en posant $P=\mathbf{P}(\mathcal{E})$,
$P'=\mathbf{P}(\mathcal{E}')$, on a
(\hyperref[II.3.5.4-fr]{3.5.4}), à un isomorphisme près,
\[
\label{II.4.1.3.2-fr}
  \mathcal{O}_{P'}(n)=\mathcal{O}_P(n)\otimes_Y\mathcal{O}_{Y'}
  \qquad\text{pour tout }n\in\mathbf{Z}.
\tag{4.1.3.2}
\]
\end{env}

